\subsection{Typography and layout}

Use only sans-serif typeface throughout.

Font sizes should follow as such:
\begin{itemize}
    \item Document title: 28
    \item Major headings (Heading 1): 20
    \item Minor Headings (Heading 2): 16
    \item Sub Headings (Heading 3): 14
    \item Body text: 12
    \item Captions, footnotes, and metadata: no smaller than 9 pt when printed
\end{itemize}

Use left-aligned text with a ragged right edge; avoid fully justified text.

Avoid relying on italics, underlining, colour, or capitalisation alone to convey meaning.

Use bold only for emphasis, labels, and headings. Not for large blocks of text.

\donot~use all capital body text.

Keep headings with the following paragraph.

Avoid leaving headings alone at the bottom of a page.

Use consistent spacing before and after headings, lists, tables, and figures.

Avoid decorative elements that consume space or reduce print clarity.

\subsection{Page and print standards}

margins should follow as:
\begin{itemize}
    \item Left: 15mm
    \item Right: 15mm
    \item Top: 15mm
    \item Bottom: 15mm
    \item Binding offset: 10 mm
\end{itemize}

Ensure the binding offset is applied to the inner edge of facing pages, not always to the left edge.

Use a page size of A4.

Use black text on a white background for normal content.

Ensure every page remains understandable when printed in grayscale.

Keep tables and figures within the printable area.

\donot~place essential information in headers or footers.

Use page numbers in the bottom right corner, except where a document control system requires a different position.

Include the document identifier and revision number in the header or footer.

\subsection{Writing conventions}

Write using Australian english spelling.

Write using direct, imperative language: “Close the valve” not “The valve should be closed.”

Use consistent terminology.

Use one action per numbered step where possible.

Number steps in the order they must be performed.

Use substeps for dependent actions, such as 1.1, 1.2, and 1.3.

State quantities, units, tolerances, temperatures, timings, and acceptance criteria explicitly.

Use SI units consistently, with a space between the number and unit: 25 °C, 10 mL, and 5 min.

Use a leading zero for values below one: 0.5 mL

Use a zero after a decimal point where appropriate: 5.0 mL

Define the required precision and rounding rules for measurements.

Use one date format throughout, preferably yyyy-mm-dd

Define abbreviations at first use and maintain an approved abbreviation list.

Use the same name for the same item throughout the document.

Avoid ambiguous terms such as “regularly”, “as needed”, “soon” or “a small amount” unless they are defined.

Avoid unexplained symbols and specialist abbreviations.

Specify what to do when a result is outside limits or a step cannot be completed.

State required prerequisites, training, equipment, materials, and personal protective equipment before the procedure begins.

Include a clear “stop and escalate” condition for unsafe or unexpected situations.

\subsection{Warnings, hazards, and precautions}

Warnings \must~appear immediately before the step or condition to which they apply.

A warning \must~state the hazard, consequence, and required action.

Red may be used as a secondary visual cue, but meaning \must~remain clear when printed in grayscale.

Safety information \must~not be conveyed by colour alone.

Warning symbols \should~be used only where their meaning is defined and consistent.

Warnings \shouldnot~be buried in long paragraphs or placed only at the end of a procedure.

Major warnings \may~be placed in a center aligned box with an all red boarder and no fill.

\subsection{Job safety analysis and task risk assessments}

See~\ref{tab:JSAExample} for a blank example JSA.

A job safety analysis (JSA) \must~be completed for tasks that are new, non-routine, changed, potentially hazardous, or insufficiently covered by an existing approved procedure.

A JSA \must\:
\begin{itemize}
    \item identify the task and its location;
    \item identify the people involved and their roles;
    \item divide the task into logical steps;
    \item identify hazards associated with each step;
    \item describe possible consequences;
    \item identify existing controls;
    \item assess the initial risk;
    \item specify additional controls where required;
    \item assess the residual risk after applying the controls;
    \item identify required personal protective equipment;
    \item identify emergency, spill, or abnormal-condition arrangements; and
    \item record the names, signatures, and dates of the people who completed and approved the assessment.
\end{itemize}

Controls \must~be selected using the following order of preference:
\begin{enumerate}
    \item elimination of the hazard;
    \item substitution with a less hazardous material, method, or process;
    \item engineering controls;
    \item administrative controls and procedural controls; and
    \item personal protective equipment.
\end{enumerate}

Personal protective equipment \mustnot~be used as the only control where a more effective control is reasonably practicable.

A JSA \must~be reviewed when the task, equipment, materials, location, personnel, or conditions change. It \must~also be reviewed after an incident, near miss, control failure, or significant unexpected result.

\subsection{Risk assessment matrix}

Risk \must{} be assessed using the likelihood and consequence scales in Tables~\ref{tab:likelihood} and~\ref{tab:consequence}. The risk score \must{} be calculated by multiplying the likelihood score by the consequence score.

\subsubsection{Likelihood scale}

The likelihood scale is defined as follows.

\begin{table}[H]
\centering
\caption{Likelihood scale}
\label{tab:likelihood}
\begin{tabular}{p{0.12\textwidth} p{0.25\textwidth} p{0.48\textwidth}}
\hline
\textbf{Score} & \textbf{Rating} & \textbf{Description} \\
\hline
1 & Rare & The event is unlikely to occur during the task or project. \\
2 & Unlikely & The event could occur, but is not expected under normal conditions. \\
3 & Possible & The event may occur at some time during the task or project. \\
4 & Likely & The event is expected to occur in some circumstances. \\
5 & Almost certain & The event is expected to occur frequently or repeatedly. \\
\hline
\end{tabular}
\end{table}

\subsubsection{Consequence scale}

The consequence scale is defined as follows.

\begin{table}[H]
\centering
\caption{Consequence scale}
\label{tab:consequence}
\begin{tabular}{p{0.12\textwidth} p{0.25\textwidth} p{0.48\textwidth}}
\hline
\textbf{Score} & \textbf{Rating} & \textbf{Description} \\
\hline
1 & Insignificant & No injury, negligible damage, and no significant disruption. \\
2 & Minor & Minor injury or illness, minor damage, or short-term disruption. \\
3 & Moderate & Medical treatment, significant damage, or disruption requiring management. \\
4 & Major & Serious injury or illness, major damage, or extended disruption. \\
5 & Catastrophic & Fatality, permanent disabling injury, multiple serious injuries, or major loss. \\
\hline
\end{tabular}
\end{table}

\subsubsection{The Matrix}

The consequence scale and likelihood scale then gives the matrix.

\begin{table}[H]
\centering
\caption{Risk matrix}
\label{tab:risk-matrix}
\renewcommand{\arraystretch}{1.5}
\begin{tabular}{c|ccccc}
\textbf{Risk Matrix}
&
\textbf{1 Insignificant}
&
\textbf{2 Minor}
&
\textbf{3 Moderate}
&
\textbf{4 Major}
&
\textbf{5 Catastrophic}
\\
\hline
\textbf{1 Rare} & 1 & 2 & 3 & 4 & 5 \\
\textbf{2 Unlikely} & 2 & 4 & 6 & 8 & 10 \\
\textbf{3 Possible} & 3 & 6 & 9 & 12 & 15 \\
\textbf{4 Likely} & 4 & 8 & 12 & 16 & 20 \\
\textbf{5 Almost certain} & 5 & 10 & 15 & 20 & 25 \\
\hline
\multicolumn{6}{c}{\textbf{Consequence score}} \\
\end{tabular}
\end{table}


\begin{description}
    \item[1--4: Low] The task may proceed with routine controls and normal supervision.

    \item[5--9: Moderate] Additional controls \should{} be considered. The responsible person \must{} confirm that the risk is adequately controlled before work begins.

    \item[10--16: High] Additional controls \must{} be implemented before the task begins. Approval by the responsible supervisor \must{} be obtained.

    \item[17--25: Extreme] The task \mustnot{} begin or continue until the risk has been reduced to an acceptable level and formally approved by the appropriate authority.
\end{description}

The risk assessment \must{} record both the initial risk and the residual risk after controls have been applied. A lower residual score \must{} not be claimed unless the additional controls are actually implemented and verified.

Where a risk cannot be reduced to an acceptable level, the task \must{} be stopped and escalated to the responsible supervisor or safety authority.


\subsection{Take 5 assessments}

See~\ref{tab:take5Example} for an blank take 5 form.

A Take 5 assessment \must~be completed immediately before starting a task when a formal JSA is not required or when the conditions have changed since the JSA or procedure was completed.

The Take 5 assessment \must~require the worker to:
\begin{enumerate}
    \item stop before starting the task;
    \item look for hazards and changes in the work area;
    \item assess the possible risk;
    \item control the hazards; and
    \item confirm that it is safe to proceed.
\end{enumerate}

A Take 5 assessment \must~be repeated when the worker, task, equipment, materials, location, or conditions change. A worker \must~stop and obtain assistance if the risk cannot be controlled adequately.

%SLAM assessments \must{} require the worker to:
%\begin{description}
%    \item[Stop] Stop and prepare before beginning the task.
%    \item[Look] Look for hazards, changes, interactions, and unsafe conditions.
%    \item[Assess] Assess the likelihood and consequence of harm.
%    \item[Manage] Apply controls, confirm that they are effective, and proceed only when the task is adequately controlled.
%\end{description}

%Take 5 and SLAM records \must{} identify the task, date, location, person conducting the assessment, hazards identified, controls applied, and any required escalation or follow-up action.

Take 5 records \must~identify the task, date, location, person conducting the assessment, hazards identified, controls applied, and any required escalation or follow-up action.

\subsection{Tables, figures, and equations}

Give every table and figure a number and descriptive title.

Refer to each table or figure in the surrounding text.

Keep table headings visible if a table continues onto another page.

\donot~use colour as the only way to distinguish table categories.

Align numbers consistently, especially decimal values.

Use sufficient cell padding and avoid overly dense tables.

Use diagrams only when they communicate something more clearly than text.

Ensure all figures have alternative text or a descriptive caption.

Avoid screenshots of text that could instead be searchable, selectable text.

\subsection{Forms, records, and examples}

Documents that explain or demonstrate the use of a form, record sheet, checklist, or template \must~include a clean, blank version of that form in Appendix A.

The example used in the main body \may~contain sample entries, annotations, instructions, or highlighting. The clean version in Appendix A \must\:
\begin{itemize}
    \item Contain no example data;
    \item Contain no instructional annotations unless they are part of the form itself;
    \item Retain all required fields, headings, signatures, dates, and document-control information;
    \item Be suitable for printing or electronic completion; and
    \item Use the same field names and layout as the example in the main body.
\end{itemize}

The example form \may~be of a different page size and orientation to the rest of the document.

Blank form pages do not require the following:
\begin{itemize}
    \item Page numbering
    \item Headers
    \item Footers
\end{itemize}

If more than one form is included, the clean forms \must~be labelled.

All forms \must~be placed in Appendix A.

A clean form \must~be provided even when the example is completed electronically. Blank fields \must~be sufficiently large for the expected amount of information.

Forms \must~not rely solely on colour, shading, or annotations to communicate their purpose.

\subsection{Accessibility and digital use}

Use real heading levels rather than manually enlarged or bold text.

Use selectable text, not scanned pages, whenever possible.

Add descriptive alternative text to meaningful images.

Use descriptive link text rather than “click here.”

Ensure hyperlinks are distinguishable without relying only on colour.

Use a logical reading order.

Make the document usable at 200\% zoom without losing essential information.

Use sufficient colour contrast for any permitted coloured content.

Include bookmarks and searchable metadata in the final document where feasible.

\subsection{Document control}

Every controlled SOP \should~include where applicable:
\begin{itemize}
    \item Document title
    \item Document identifier
    \item Version or revision number
    \item Effective date
    \item Review date
    \item Author
    \item Technical reviewer
    \item Approver
    \item Approval date
    \item Superseded document information
    \item Page number in the form “Page X of Y,” if practical
    \item A revision-history table
    \item A statement identifying whether the copy is controlled or uncontrolled
\end{itemize}